\documentclass[10pt,a4paper]{amsart}
\usepackage[margin=19mm]{geometry}
\usepackage{amsmath,amssymb,mathtools}
\usepackage[scr=rsfs,cal=euler]{mathalfa}
\usepackage{xfrac}
\usepackage[unicode,hidelinks,bookmarks=false]{hyperref}
\newcommand{\ie}{i.e.\ }
\newcommand{\cf}{cf.\ }
\newcommand{\resp}{respectively\ }
\newcommand{\Subsection}{\subsection}
\providecommand{\nobreakdash}{\nobreak\hbox{-}}
\setlength{\emergencystretch}{2em}
\multlinegap0pt
\usepackage{amssymb}
\usepackage{mathtools}
\usepackage{tcolorbox}
\newtheorem{theorem}{Theorem}
\newcommand{\LC}{\Delta_{\mathrm{LC}}}
\newcommand{\Pbb}{\mathbb P}
\newcommand{\Poi}{\mathrm{Poi}}
\newcommand{\R}{\mathbb R}
\newcommand{\TV}{\mathrm{TV}}
\title{Universal Shuffle Asymptotics, Part III:\\ Dominant-Block Quotient Geometry and Hybrid Gaussian--Compound-Poisson Limits\\ in Finite-Alphabet Shuffle Privacy}
\author{Alex Shvets}
\date{Source: arXiv:2603.13407v1 (CC BY 4.0)}
\begin{document}
\maketitle

\noindent\emph{Source and licence.} Universal Shuffle Asymptotics, Part III:\\ Dominant-Block Quotient Geometry and Hybrid Gaussian--Compound-Poisson Limits\\ in Finite-Alphabet Shuffle Privacy. Version arXiv:2603.13407v1 (CC BY 4.0), distributed by arXiv under the Creative Commons Attribution 4.0 International licence, http://creativecommons.org/licenses/by/4.0/, as recorded in the arXiv licence metadata for this exact version and shown on the abstract page https://arxiv.org/abs/2603.13407v1. Full licence notice: \texttt{licenses/arxiv-2603.13407-CC-BY-4.0.txt}.\par\smallskip

\noindent\textit{Editorial scope.} Counted unit: the article's theorem thm:boundary-BE-poisson (source0.tex lines 1851--1872). The article's complete original proof of it is delivered with it (source0.tex lines 1874--2026, one \texttt{proof} environment of 153 lines). No mathematics is written, repaired or re-derived: the statement and the proof are the article's own lines, in the article's own notation, and no retained line is substituted or re-flowed.  No further source-local context is needed: every cross-reference of every retained line resolves inside the two delivered spans.  Nothing was omitted from any retained span, so the package carries no omission record.  No retained line contains a citation, so the wrapper carries no bibliography and no bibliography entry.  No retained line contains a figure environment, an image-inclusion command or drawing code, so no image is delivered and the package declares no asset dependency.  Editorial formatting only, in addition: an AMS \texttt{amsart} preamble that restates the article's own declarations and macros used by the retained lines, namely the environments \texttt{theorem} on separate counters, together with the standard packages those lines need.  The retained environments print in this document as Theorem 1 in order of appearance, whereas the article prints the same items with the numbering of its own full document.  The complete native source of the article as distributed in the arXiv e-print for this exact version is bundled unchanged as \texttt{sources/196330/source0.tex}.
\begin{theorem}[Boundary Berry--Esseen for the Poisson-shift experiment]
\label{thm:boundary-BE-poisson}
There exist constants $c_0>0$ and $C<\infty$ such that for every $c\in(0,c_0]$,
\[
\LC\bigl(\mathcal E_c^{\mathrm{Poi}},\mathcal E_c^{\mathrm G}\bigr)\le Cc.
\]
More concretely, after the rescaling
\[
Z_c:=c\bigl(J_c-c^{-2}\bigr),\qquad J_c\sim\Poi(c^{-2}),
\]
the experiment $\mathcal E_c^{\mathrm{Poi}}$ is equivalent to the lattice experiment
\[
\widehat{\mathcal E}_c^{\mathrm{Poi}}:=\bigl(\mathcal L(Z_c),\ \mathcal L(Z_c+c)\bigr),
\]
and there exist Markov kernels $K_c$ and $L_c$ such that
\[
\TV\bigl(K_c\mathcal L(Z_c),N(0,1)\bigr)+\TV\bigl(K_c\mathcal L(Z_c+c),N(c,1)\bigr)\le Cc,
\]
\[
\TV\bigl(L_cN(0,1),\mathcal L(Z_c)\bigr)+\TV\bigl(L_cN(c,1),\mathcal L(Z_c+c)\bigr)\le Cc.
\]
\end{theorem}

\begin{proof}
Set $\lambda:=c^{-2}$ and let
\[
z_{k,c}:=c(k-\lambda),\qquad k\in\mathbb Z.
\]
The mapping $k\mapsto z_{k,c}$ is bijective, so $\mathcal E_c^{\mathrm{Poi}}$ is exactly equivalent to the lattice
experiment $\widehat{\mathcal E}_c^{\mathrm{Poi}}$.

Let $\varphi$ be the standard normal density and write $I_{k,c}:=[z_{k,c}-c/2,z_{k,c}+c/2)$.
We claim that
\begin{equation}
\Pbb(Z_c=z_{k,c})
=
c\varphi(z_{k,c})+O\!\left(c^2(1+|z_{k,c}|)^{-3}\right)
\label{eq:poisson-local-weighted}
\end{equation}
uniformly in $k$ as $c\downarrow 0$.

Indeed, let $p_{k,\lambda}:=\Pbb(\Poi(\lambda)=k)=e^{-\lambda}\lambda^k/k!$ with $\lambda=c^{-2}$ and
$u:=z_{k,c}=c(k-\lambda)$. For $k<0$ the left-hand side is zero, so the claimed bound is immediate; hence assume
$k\ge 0$.

If $|u|\le c^{-1/3}$, then $k=\lambda+u/c=\lambda(1+uc)$ with $|uc|\le c^{2/3}$. Stirling's formula gives
\[
p_{k,\lambda}
=
\frac{1}{\sqrt{2\pi k}}
\exp\!\bigl(-\lambda+k-k\log(k/\lambda)\bigr)\bigl(1+O(k^{-1})\bigr).
\]
Since
\[
(1+t)\log(1+t)-t=\frac{t^2}{2}+O(|t|^3)
\qquad (|t|\le 1/2),
\]
we obtain
\[
-\lambda+k-k\log(k/\lambda)
=
-\lambda\bigl((1+uc)\log(1+uc)-uc\bigr)
=
-\frac{u^2}{2}+O\!\bigl(c|u|^3\bigr),
\]
while
\[
k^{-1/2}=c\bigl(1+O(c(1+|u|))\bigr).
\]
Hence, uniformly on $|u|\le c^{-1/3}$,
\[
p_{k,\lambda}
=
c\,\varphi(u)
\exp\!\bigl(O(c|u|^3)\bigr)
\bigl(1+O(c(1+|u|))\bigr).
\]
We now split this central range into two subranges.

If $|u|\le c^{-1/4}$, then $c|u|^3=O(c^{1/4})$, so the exponential may be linearized and
\[
p_{k,\lambda}
=
c\,\varphi(u)\Bigl(1+O\!\bigl(c(1+|u|^3)\bigr)\Bigr)
=
c\,\varphi(u)+O\!\left(c^2(1+|u|)^{-3}\right),
\]
because $\sup_{u\in\R}\varphi(u)(1+|u|)^6<\infty$.

If instead $c^{-1/4}<|u|\le c^{-1/3}$, then $u^2\ge c^{-1/2}$ and $c|u|^3\le c^{1/4}$. For all sufficiently
small $c$ this gives
\[
p_{k,\lambda}\le Cc\,e^{-u^2/2+c^{1/4}}\le Cc\,e^{-u^2/4}\le Cc\,e^{-c^{-1/2}/4},
\]
while also
\[
c\,\varphi(u)\le Cc\,e^{-u^2/2}\le Cc\,e^{-c^{-1/2}/2}.
\]
Both quantities are super-exponentially small in $c^{-1/2}$, hence in particular they are
$O\!\left(c^2(1+|u|)^{-3}\right)$.

Finally, if $|u|>c^{-1/3}$, then $|k-\lambda|=|u|/c\ge c^{-4/3}=\lambda^{2/3}$. If $k\ge \lambda+\lambda^{2/3}$,
then $p_{k,\lambda}\le \Pbb(\Poi(\lambda)\ge k)$, and the standard one-sided Chernoff bound gives
\[
p_{k,\lambda}\le e^{-c_1\lambda^{1/3}}=e^{-c_1c^{-2/3}}
\]
for some absolute constant $c_1>0$. If $0\le k\le \lambda-\lambda^{2/3}$, then similarly
\[
p_{k,\lambda}\le \Pbb(\Poi(\lambda)\le k)\le e^{-c_1c^{-2/3}}.
\]
Moreover,
\[
c\,\varphi(u)\le c\,e^{-u^2/2}\le c\,e^{-c^{-2/3}/2}.
\]
Again both bounds are $O\!\left(c^2(1+|u|)^{-3}\right)$. This proves
\eqref{eq:poisson-local-weighted}.

Since $|\varphi'(x)|\le C(1+|x|)^{-3}$ for all $x$, integration over the cell $I_{k,c}$ yields
\[
\int_{I_{k,c}}\varphi(x)\,dx
=
c\varphi(z_{k,c})+O\!\left(c^2(1+|z_{k,c}|)^{-3}\right),
\]
and therefore
\begin{equation}
\left|\Pbb(Z_c=z_{k,c})-\int_{I_{k,c}}\varphi(x)\,dx\right|
\le \frac{C_1c^2}{(1+|z_{k,c}|)^3}.
\label{eq:poisson-cell-bound}
\end{equation}

For the deficiency from Poisson to Gaussian, let $K_c$ be the uniform-jitter kernel
\[
K_c(x,\cdot):=\mathrm{Unif}[x-c/2,x+c/2].
\]
The law $K_c\mathcal L(Z_c)$ has density
\[
f_c(x)=\frac{1}{c}\Pbb(Z_c=z_{k,c})\qquad \text{for }x\in I_{k,c}.
\]
Let $\bar\varphi_{k,c}:=c^{-1}\int_{I_{k,c}}\varphi(u)\,du$. Then
\[
\int_{I_{k,c}}|f_c(x)-\varphi(x)|\,dx
\le \left|\Pbb(Z_c=z_{k,c})-\int_{I_{k,c}}\varphi(u)\,du\right|
+\int_{I_{k,c}}|\bar\varphi_{k,c}-\varphi(x)|\,dx.
\]
The first term is bounded by \eqref{eq:poisson-cell-bound}, and the second is at most
$Cc^2(1+|z_{k,c}|)^{-3}$ by the derivative bound on $\varphi$. Hence
\[
\int_{I_{k,c}}|f_c(x)-\varphi(x)|\,dx\le \frac{C_2c^2}{(1+|z_{k,c}|)^3}.
\]
Summing over $k$ and using that $z_{k,c}$ runs over a lattice of mesh $c$,
\[
\TV\bigl(K_c\mathcal L(Z_c),N(0,1)\bigr)
\le C_2c^2\sum_{k\in\mathbb Z}(1+|z_{k,c}|)^{-3}
\le C_3c.
\]
By translation invariance of both the jitter kernel and the Gaussian family,
\[
\TV\bigl(K_c\mathcal L(Z_c+c),N(c,1)\bigr)
=\TV\bigl(K_c\mathcal L(Z_c),N(0,1)\bigr)
\le C_3c.
\]

For the reverse deficiency, let $L_c$ be the rounding kernel that maps $x\in I_{k,c}$ to $z_{k,c}$. Then
$L_cN(0,1)$ is the lattice law assigning mass $\int_{I_{k,c}}\varphi(x)\,dx$ to $z_{k,c}$, so by
\eqref{eq:poisson-cell-bound},
\[
\TV\bigl(L_cN(0,1),\mathcal L(Z_c)\bigr)
\le \frac12\sum_{k\in\mathbb Z}\frac{C_1c^2}{(1+|z_{k,c}|)^3}
\le C_4c.
\]
Again translation invariance yields
\[
\TV\bigl(L_cN(c,1),\mathcal L(Z_c+c)\bigr)\le C_4c.
\]
Taking $C:=\max\{C_3,C_4\}$ proves the theorem.
\end{proof}

\end{document}
